\documentclass[11pt,a4paper]{article}

\usepackage{fontspec}
\setmainfont{DejaVu Sans}
\setsansfont{DejaVu Sans}
\setmonofont{DejaVu Sans Mono}
\usepackage[english]{babel}
\usepackage[a4paper,margin=2.1cm]{geometry}
\usepackage{xcolor}
\usepackage{hyperref}
\usepackage{xurl}
\usepackage{booktabs}
\usepackage{longtable}
\usepackage{array}
\usepackage{enumitem}
\usepackage{fancyhdr}
\usepackage{titlesec}
\usepackage{listings}
\usepackage{caption}
\usepackage{lastpage}

\definecolor{primary}{HTML}{1F4E79}
\definecolor{darkgray}{HTML}{333333}
\definecolor{lightgray}{HTML}{F3F6F8}
\definecolor{codegray}{HTML}{F7F7F7}

\hypersetup{
    colorlinks=true,
    linkcolor=primary,
    urlcolor=primary,
    pdftitle={Technical recruiter dossier - C++ Space Telemetry Ground Station},
    pdfauthor={Guillaume Lemonnier}
}

\pagestyle{fancy}
\fancyhf{}
\lhead{C++ Space Telemetry Ground Station}
\rhead{Technical recruiter dossier}
\cfoot{\thepage\ / \pageref{LastPage}}

\titleformat{\section}{\Large\bfseries\color{primary}}{\thesection}{0.7em}{}
\titleformat{\subsection}{\large\bfseries\color{darkgray}}{\thesubsection}{0.7em}{}
\titleformat{\subsubsection}{\normalsize\bfseries\color{darkgray}}{\thesubsubsection}{0.7em}{}

\lstset{
    basicstyle=\ttfamily\small,
    backgroundcolor=\color{codegray},
    frame=single,
    rulecolor=\color{lightgray},
    breaklines=true,
    columns=fullflexible,
    keepspaces=true,
    showstringspaces=false,
    tabsize=2
}

\newcommand{\projectname}{C++ Space Telemetry Ground Station}
\newcommand{\code}[1]{\texttt{#1}}
\newcommand{\decision}[2]{\noindent\fcolorbox{primary}{lightgray}{\begin{minipage}{0.96\linewidth}\textbf{#1}\par\vspace{0.2em}#2\end{minipage}}\par\vspace{0.8em}}

\begin{document}
\sloppy
\emergencystretch=3em

\begin{titlepage}
    \centering
    \vspace*{1.8cm}
    {\Huge\bfseries\color{primary} \projectname \par}
    \vspace{0.5cm}
    {\Large Technical dossier for software recruiters\par}
    \vspace{1.2cm}
    {\large English version\par}
    \vspace{1.0cm}
    \rule{0.82\linewidth}{0.6pt}\par
    \vspace{0.6cm}
    {\large C++20 portfolio project - Linux - telemetry - networking - multithreading - CMake - testing - CI/CD\par}
    \vspace{0.6cm}
    \rule{0.82\linewidth}{0.6pt}\par
    \vfill
    \begin{flushleft}
    \textbf{Document purpose}\par
    Present the project to a recruiter or software engineer by explaining the implementation goals, architecture, development choices, code quality practices and skills demonstrated by the codebase.\par
    \vspace{0.7cm}
    \textbf{Author}\par
    Guillaume Lemonnier\par
    \vspace{0.3cm}
    \textbf{Date}\par
    July 2026\par
    \end{flushleft}
\end{titlepage}

\tableofcontents
\newpage

\section{Executive summary}

\projectname{} is a Linux-based C++20 project that simulates a ground station able to receive, decode, validate, replay and export satellite-like or embedded-device telemetry frames. It was designed as a skills demonstration for roles involving industrial software, systems programming, telemetry, networking, modern C++ and constrained environments.

The project covers an end-to-end chain: telemetry frame generation, UDP/TCP reception, TCP stream reconstruction, strict binary parsing, CRC-32 validation, a multithreaded processing pipeline, replay files, CSV/JSON export, logging, automated tests, GitHub Actions CI, reproducible benchmarking and operational degraded mode.

\decision{Recruiter positioning}{This project is not intended to be a certified space product. It demonstrates the ability to build a small but structured, testable, robust and explainable C++ system with development choices close to industrial, defense, space or embedded-software concerns.}

\section{Implementation intent}

The project has three main goals.

\subsection{Demonstrate practical modern C++}

The goal was not to write a small syntax exercise, but a project where C++ is used to structure a real application pipeline: domain models, binary encoding, error handling, Linux I/O, thread-safe queues, tests and command-line tools.

The project exercises:

\begin{itemize}[nosep]
    \item C++20, RAII, strong domain types and modular design;
    \item \code{std::variant} to distinguish a valid frame from a parsing error;
    \item closeable thread-safe queues to decouple reception, decoding and writing;
    \item strict CMake builds, warnings and automated testing.
\end{itemize}

\subsection{Move the project toward an industrial context}

Telemetry was chosen deliberately because it introduces systems-oriented topics: binary protocols, integrity checks, network loss, data corruption, replay, state monitoring and performance measurement.

The project therefore includes:

\begin{itemize}[nosep]
    \item a strict binary protocol;
    \item CRC validation before trusting telemetry data;
    \item loss and corruption simulation;
    \item degraded mode when the incoming stream becomes too unstable;
    \item a benchmark to measure pipeline behavior.
\end{itemize}

\subsection{Produce an interview-ready project}

The codebase was organized to be easy to present: clear structure, technical documentation, English README, reproducible commands and separation between domain code, networking infrastructure, replay, tests and benchmarking.

\section{Delivered functional scope}

\begin{longtable}{p{0.18\linewidth} p{0.28\linewidth} p{0.44\linewidth}}
\toprule
\textbf{Version} & \textbf{Goal} & \textbf{Implemented features} \\
\midrule
V1 & Core pipeline & Telemetry frame simulator, UDP receiver, C++ parser, CRC-32 validation, logs, unit tests, English README. \\
\addlinespace
V2 & Industrialization & Multithreaded pipeline, STGF replay files, CSV and JSON export, reproducible benchmark, CTest, GitHub Actions. \\
\addlinespace
V3 & Operational robustness & Stricter binary protocol, network loss and corruption simulation, degraded mode, architecture documentation, performance report. \\
\bottomrule
\end{longtable}

\section{Overall architecture}

The architecture is deliberately split into simple components. The ground station can receive frames from the network or from a replay file. Raw frames are pushed into a queue, decoded by one or more workers, monitored by a health component and then exported.

\begin{lstlisting}
                  +-------------------------+
                  | stgs_satellite_simulator|
                  |  UDP/TCP or STGF file   |
                  +------------+------------+
                               |
                               v
+------------------+    +------+-------+    +----------------+    +---------------+
| Network receiver | -> | Raw queue    | -> | Decoder workers| -> | Decoded queue |
| or replay reader |    | ByteVector   |    | CRC + parsing  |    | TelemetryFrame|
+------------------+    +--------------+    +------+---------+    +-------+-------+
                                                   |                      |
                                                   v                      v
                                          +--------+---------+      +------+------+
                                          | Health monitor   |      | CSV/JSON    |
                                          | nominal/degraded |      | writer      |
                                          +------------------+      +-------------+
\end{lstlisting}

\subsection{Code organization}

\begin{longtable}{>{\raggedright\arraybackslash}p{0.44\linewidth} >{\raggedright\arraybackslash}p{0.46\linewidth}}
\toprule
\textbf{File or folder} & \textbf{Responsibility} \\
\midrule
\path|apps/ground_station.cpp| & Main application: CLI, UDP/TCP reception, replay, workers, exports and degraded mode. \\
\path|apps/satellite_simulator.cpp| & Telemetry generator: network output, STGF file output, loss, corruption and simulated rate. \\
\path|include/stgs/TelemetryFrame.hpp| & Domain model for a decoded telemetry frame. \\
\path|include/stgs/FrameCodec.hpp|\newline\path|src/FrameCodec.cpp| & Encoding, decoding, validation and frame extraction from a TCP stream. \\
\path|include/stgs/NetworkServer.hpp|\newline\path|src/NetworkServer.cpp| & UDP/TCP reception on Linux using POSIX sockets and \code{poll()}. \\
\path|include/stgs/BlockingQueue.hpp| & Closeable producer/consumer queue for the multithreaded pipeline. \\
\path|include/stgs/StationHealth.hpp|\newline\path|src/StationHealth.cpp| & Operational health monitoring and nominal/degraded transitions. \\
\path|include/stgs/Replay.hpp|\newline\path|src/Replay.cpp| & Reading and writing the STGF replay format. \\
\path|benchmarks/benchmark_decode.cpp| & Reproducible benchmark for the decode pipeline. \\
\path|tests/test_main.cpp| & Unit tests and behavioral tests. \\
\path|.github/workflows/ci.yml| & Build, tests, smoke replay, JSON export validation and benchmark smoke. \\
\bottomrule
\end{longtable}

\section{Development choices}

\subsection{Why C++20}

C++20 keeps the project close to systems and industrial software roles while providing a more expressive and safer style than legacy C++. The project favors explicit design: domain types, isolated functions, testable components and very few external dependencies.

\decision{Why few dependencies?}{The project is intentionally lightweight. The goal is to show mastery of C++/Linux fundamentals rather than hide complexity behind frameworks. This also keeps the build simple in CI and makes the code easier to audit.}

\subsection{Why CMake}

CMake structures the project into separate targets. The main targets are:

\begin{itemize}[nosep]
    \item \code{stgs\_core} for core components;
    \item \code{stgs\_network} for networking infrastructure;
    \item \code{stgs\_ground\_station} for the receiver application;
    \item \code{stgs\_satellite\_simulator} for the telemetry generator;
    \item \code{stgs\_tests} for CTest;
    \item \code{stgs\_benchmark\_decode} for performance measurements.
\end{itemize}

This separation allows the core library to be built independently from the applications and reused by tests.

\begin{lstlisting}[language=bash]
cmake -S . -B build -DCMAKE_BUILD_TYPE=Release
cmake --build build -j
ctest --test-dir build --output-on-failure
\end{lstlisting}

\subsection{Error model}

The parser returns a \code{std::variant<TelemetryFrame, FrameError>}. Invalid telemetry is treated as an expected data condition, not as an application exception.

This is appropriate for telemetry streams: in a noisy network, invalid data should not crash the process. It should be counted, logged and used as an input for degraded-mode decisions.

\begin{lstlisting}[language=C++]
using FrameParseResult = std::variant<TelemetryFrame, FrameError>;

FrameParseResult decodeFrame(std::span<const std::uint8_t> bytes);
\end{lstlisting}

\subsection{Strict binary protocol}

The protocol contains a magic word, a version, a satellite identifier, a timestamp, measurements, a payload length and a CRC-32. All integers are encoded in big-endian order to match network byte order.

\begin{lstlisting}
MAGIC | VERSION | SATELLITE_ID | TIMESTAMP_MS | TEMPERATURE_C |
BATTERY_PERCENT | STATUS | PAYLOAD_LEN | PAYLOAD | CRC32
\end{lstlisting}

The validation rules are strict: expected magic word, supported version, coherent length, maximum payload size, valid CRC, battery between 0 and 100 and known status.

\decision{Why a magic word?}{With TCP, the transport is a byte stream and does not preserve application frame boundaries. The STGS magic word allows the receiver to resynchronize the stream if noise or corrupted bytes appear.}

\subsection{UDP and TCP choices}

The project implements both approaches:

\begin{itemize}[nosep]
    \item in UDP mode, each datagram is treated as one complete candidate frame;
    \item in TCP mode, an extractor rebuilds frames from a byte stream.
\end{itemize}

This illustrates two networking models: simplicity and expected loss with UDP, transport reliability but application-level framing requirements with TCP.

\section{Multithreaded pipeline}

The pipeline separates responsibilities: reception, decoding, health monitoring and writing. The number of decoder workers is configurable with \code{--decoder-threads}. This design demonstrates concurrency without adding unnecessary complexity.

\begin{lstlisting}
Receiver or replay reader
    -> BlockingQueue<ByteVector>
        -> decoder worker 1
        -> decoder worker 2
        -> ...
            -> BlockingQueue<TelemetryFrame>
                -> CSV/JSON writer
\end{lstlisting}

The \code{BlockingQueue} component is closeable, which makes thread shutdown simpler. Once no more data is available, workers can exit cleanly without busy-waiting.

\section{Degraded mode}

\subsection{Intent}

Degraded mode was added deliberately to go beyond simple frame parsing. It introduces the idea of an operational ground-station state: even if the program keeps running, the telemetry stream can become poor enough to require a degraded-state signal.

\subsection{Mechanism}

\code{StationHealthMonitor} maintains a sliding window of recent samples. Each sample is either a rejected frame or a decoded frame. A decoded frame is considered critical if its status is \code{CRITICAL} or \code{SAFE\_MODE}.

The station moves from \code{NOMINAL} to \code{DEGRADED} if:

\begin{itemize}[nosep]
    \item the rejection rate reaches the configured threshold;
    \item or the number of critical frames reaches the configured threshold.
\end{itemize}

It returns to \code{NOMINAL} when the rejection rate falls below the recovery threshold and the window no longer contains a problematic critical volume.

\begin{lstlisting}
--degraded-window 100
--degraded-min-samples 20
--degraded-rejection-rate 0.10
--degraded-recovery-rate 0.03
--degraded-critical-count 3
\end{lstlisting}

\decision{Value for a recruiter}{This part demonstrates the ability to think about operational behavior, not only nominal output. The code monitors errors, makes a state decision and logs transitions.}

\section{Replay, export and observability}

\subsection{STGF replay format}

The STGF format records binary frames and replays them through the same pipeline as live network data. This matters for testing, reproducibility and incident analysis.

\begin{lstlisting}
'S' 'T' 'G' 'F' 0 0 0 1
FRAME_LENGTH:uint32_be | FRAME_BYTES
\end{lstlisting}

\subsection{CSV and JSON exports}

CSV enables simple spreadsheet-style inspection. JSON supports integration with analysis tools, dashboards or automated checks. The format can be inferred from the extension or forced with \code{--output-format}.

\begin{lstlisting}[language=bash]
./build/stgs_ground_station --replay frames.stgf --output replay.csv
./build/stgs_ground_station --replay frames.stgf --output replay.json
./build/stgs_ground_station --replay frames.stgf --output telemetry.out --output-format json
\end{lstlisting}

\subsection{Logs}

The logger is thread-safe and can write to the console or to a file. Logs help explain station behavior: startup, frame errors, health transitions and final summaries.

\section{Testing, quality and CI/CD}

\subsection{Testing strategy}

The tests cover critical behavior: known CRC vector, encode/decode round trip, rejection for bad magic, CRC mismatch, length mismatch, TCP extraction with noise, replay round trip, blocking queue closure and degraded-mode transitions.

\begin{lstlisting}[language=bash]
ctest --test-dir build --output-on-failure
\end{lstlisting}

\subsection{GitHub Actions}

The CI workflow builds in Release mode, runs CTest, generates an STGF replay file, replays to CSV, replays to JSON, validates JSON with Python and runs a benchmark smoke test. The project is therefore designed to be verified automatically, not only run locally.

\decision{Quality choice}{For a recruiter, CTest, CI, a benchmark and error-case tests are as important as the nominal feature set. They demonstrate maintainability and industrialization habits.}

\section{Benchmark and performance}

The project includes \code{stgs\_benchmark\_decode}, a tool that measures decode-pipeline throughput without network or disk I/O. It pre-generates valid frames, pushes them into the queue, decodes them with a configurable number of workers and measures frames per second.

\begin{lstlisting}[language=bash]
./build/stgs_benchmark_decode --frames 200000 --payload-size 64 --decoder-threads 4
\end{lstlisting}

The performance report documents the methodology and reference measurements. One useful observation is that multithreading does not automatically improve performance for small workloads: queue synchronization can cost more than the parsing work itself. This is a valuable interview point because it shows that concurrency choices should be measured.

\section{Skills demonstrated}

\begin{longtable}{p{0.30\linewidth} p{0.60\linewidth}}
\toprule
\textbf{Skill} & \textbf{Evidence in the project} \\
\midrule
Modern C++ & C++20, domain types, \code{std::variant}, RAII and library separation. \\
Linux & POSIX sockets, \code{poll()}, CLI tools, files, Linux build and execution. \\
Networking & UDP/TCP reception, TCP framing, loss and corruption simulation. \\
Binary protocol & Magic word, version, lengths, big-endian fields, CRC and variable payload. \\
Multithreading & Thread-safe queues, decoder workers, separate writer and clean shutdown. \\
Performance & Reproducible benchmark, frames/second measurements, payload/thread analysis. \\
CMake & Separate targets, build options for tests/benchmarks, CTest. \\
Testing & Unit tests and error-case tests, CTest and CI integration. \\
CI/CD & GitHub Actions with build, tests, replay, JSON validation and benchmark smoke. \\
Documentation & English README, architecture documentation and performance report. \\
Robustness & Invalid-frame rejection, health monitoring, degraded mode and logs. \\
\bottomrule
\end{longtable}

\section{Known limitations}

The project is intentionally educational and portfolio-oriented. It does not claim to meet the requirements of a certified space system. Its main limitations are:

\begin{itemize}[nosep]
    \item telemetry-inspired protocol, but not CCSDS-compatible;
    \item no real hardware or driver layer;
    \item no long-term persistence or file rotation;
    \item no Prometheus metrics or web dashboard;
    \item no fuzzing or full network load testing;
    \item no system package or systemd deployment unit.
\end{itemize}

These limitations are useful to mention in interviews because they show the ability to distinguish a portfolio project from production-critical software.

\section{Possible extensions}

The most relevant next improvements would be:

\begin{itemize}[nosep]
    \item add a CCSDS-inspired header with a sequence counter;
    \item estimate UDP loss from frame sequence counters;
    \item add a WebSocket dashboard or Prometheus metrics endpoint;
    \item add parser fuzzing tests;
    \item create a Linux service mode with systemd;
    \item add an embedded-device or hardware-sensor simulator;
    \item add more realistic network load scenarios;
    \item generate Doxygen documentation for the C++ core.
\end{itemize}

\section{Interview pitch}

\begin{quote}
I developed this project to demonstrate my progress in modern systems-oriented C++. The goal was to build a mini telemetry ground station able to receive UDP/TCP frames or replay files, decode them through a strict binary protocol, validate integrity with CRC, export data, measure performance and monitor operational state through degraded mode. The project primarily demonstrates my ability to structure testable C++ code, document architecture, build with CMake, verify through CI and explain technical trade-offs clearly.
\end{quote}

\section{Demo commands}

\subsection{Build and tests}

\begin{lstlisting}[language=bash]
cmake -S . -B build -DCMAKE_BUILD_TYPE=Release
cmake --build build -j
ctest --test-dir build --output-on-failure
\end{lstlisting}

\subsection{Replay to CSV and JSON}

\begin{lstlisting}[language=bash]
./build/stgs_satellite_simulator --output-file frames.stgf --count 1000 --rate 0
./build/stgs_ground_station --replay frames.stgf --output replay.csv
./build/stgs_ground_station --replay frames.stgf --output replay.json
\end{lstlisting}

\subsection{UDP simulation with loss and corruption}

\begin{lstlisting}[language=bash]
./build/stgs_ground_station --udp --port 9000 --output telemetry.json --log ground.log
./build/stgs_satellite_simulator --udp --host 127.0.0.1 --port 9000 \
  --count 5000 --rate 2000 --loss 0.02 --corrupt 0.01
\end{lstlisting}

\subsection{Benchmark}

\begin{lstlisting}[language=bash]
./build/stgs_benchmark_decode --frames 200000 --payload-size 64 --decoder-threads 4
./benchmarks/run_benchmark.sh
\end{lstlisting}

\section{Conclusion}

\projectname{} is a recruiter-oriented portfolio project that demonstrates solid foundations in modern C++, Linux, networking, binary protocols, multithreading, CMake, testing, CI/CD, robustness and technical documentation. Its main value is the complete engineering approach: define a clear technical goal, design a modular architecture, implement features, test error cases, measure performance and document decisions.

For industrial, space, defense, systems or telemetry software roles, this project is a concrete proof of structured learning and of the ability to produce explainable, maintainable and verifiable code.

\end{document}
